\documentclass[11pt]{article}
\usepackage[a4paper,margin=1in]{geometry}
\usepackage{amsmath,amssymb,amsthm,mathtools}
\usepackage{microtype}
\usepackage{enumitem}
\usepackage{xcolor}
\usepackage[colorlinks=true,linkcolor=blue!55!black,citecolor=green!40!black,
            urlcolor=blue!60!black]{hyperref}

\newtheorem{theorem}{Theorem}[section]
\newtheorem{lemma}[theorem]{Lemma}
\newtheorem{corollary}[theorem]{Corollary}
\newtheorem{proposition}[theorem]{Proposition}
\newtheorem{conjecture}[theorem]{Conjecture}
\theoremstyle{definition}
\newtheorem{definition}[theorem]{Definition}
\newtheorem{observation}[theorem]{Observation}
\theoremstyle{remark}
\newtheorem{remark}[theorem]{Remark}

\newcommand{\Q}{\mathbb Q}
\newcommand{\Z}{\mathbb Z}
\newcommand{\Zp}{\mathbb Z_{(p)}}
\newcommand{\N}{\mathbb N}
\newcommand{\ct}{\operatorname{CT}}
\newcommand{\ord}{\operatorname{ord}}
\newcommand{\e}{\mathrm e}
\newcommand{\HH}[2]{H_{#1}^{(#2)}}
\newcommand{\KH}[3]{H_{#1}^{(#2,#3)}}
\newcommand{\mix}{\partial_1\cdots\partial_r}
\newcommand{\Verified}{\textcolor{blue!65!black}{\textsc{verified}}}
\newcommand{\Open}{\textcolor{orange!75!black}{\textsc{open}}}

\title{Harmonic Jets of Hypergeometric Sums:\\
Inverse Parameter Lifting, Cyclotomic Holonomicity, and Lucas Descent}
\author{[Author]}
\date{August 2026}

\begin{document}
\maketitle

\begin{abstract}
We develop a general calculus for finite sums in which a hypergeometric kernel is
multiplied by a polynomial in ordinary or cyclotomic harmonic numbers.  The first result
is an inverse to the familiar operation of differentiating Pochhammer symbols: every
harmonic polynomial of prescribed weight is a finite linear combination of mixed
derivatives of explicit hypergeometric parameter products.  An inclusion--exclusion
product isolates a single letter $H_m^{(r)}$ without contamination by lower set
partitions.  It follows, after a finite residue-class split, that every finite proper
hypergeometric sum with polynomial cyclotomic harmonic weight is P-recursive; the proof
is constructive and reduces recurrence certification to ordinary creative telescoping.

The second result is arithmetic.  We prove a character-resolved Lucas descent theorem:
under a termwise digit factorisation and a transparent compatibility condition on the
harmonic arguments, the weight-$w$ companion satisfies
\[
 p^w B_\eta(ap+r)\equiv \eta(p)B_\eta(a)A(r)\pmod p.
\]
Different character signatures therefore form separate Frobenius eigenspaces.  Finally
we apply the lifting calculus to a sporadic recurrence whose Ap\'ery limit is one half of
Catalan's constant.  A compact two-parameter product has as its mixed derivative a new
$\chi_{-4}$-harmonic closed form for the second solution.  We prove the identity by an
exact coupled holonomic reduction: two different recurrence residuals share an order-$10$
Ore factor, and finite coefficient propagation proves that their difference is $1$.
The formula then evaluates its own limit as $G/2$ by a positive-shell concentration
argument, yields a positive error series for Catalan's constant, and gives a
character-twisted Lucas congruence on the uniform half-digit range.  A second,
Casoratian-based argument proves the unrestricted digit congruence at every prime for
which the first-solution digit block has no zero modulo $p$.  Finally, an endpoint
binomial transform decouples the companion recurrence and proves the sharp global
denominator bound $\operatorname{lcm}(1,\ldots,n)^2B_E(n)\in\mathbb Z$.  The same
transform represents $G$ as the finite part of a zero-balanced reciprocal-central-
binomial series; the full transform family has an optimally accelerated member with
convergence factor $1/3$, but also gives a sharp obstruction to obtaining irrationality
from any fixed integral binomial transform.  More generally, endpoint descent gives a
complete denominator classification for the dilation family: its normalized companion
has the $L_n^2$ bound at every index if and only if the dilation parameter is divisible
by $4$.
\end{abstract}

\section{Introduction}

Finite binomial sums with harmonic weights occur in Ap\'ery limits, congruences,
Feynman-integral calculations, and symbolic summation.  A typical object is
\begin{equation}\label{eq:typical}
 B(n)=\sum_k S(n,k)\,W(n,k),
\end{equation}
where $S$ is a proper hypergeometric term and $W$ is a polynomial in letters such as
$H_k$, $H_{n-k}^{(2)}$, or
\[
 H_m^{(q,\chi)}:=\sum_{1\le j\le m}\frac{\chi(j)}{j^q}.
\]
The standard direction is to introduce parameters into factorials and differentiate.
The reverse direction is usually performed by inspired guessing: given $W$, one searches
for a deformation that happens to produce it.  The first purpose of this paper is to make
that reverse step canonical.

Three principles emerge.
\begin{enumerate}[leftmargin=2em]
\item A harmonic letter is a \emph{connected parameter derivative}.  Boolean-lattice
  inclusion--exclusion removes every disconnected product of lower-weight letters.
\item Cyclotomic letters become ordinary hypergeometric parameters after splitting the
  relevant indices into finitely many residue classes.
\item Under Lucas factorisation, harmonic weight is a Frobenius degree: a letter of
  weight $q$ contributes $p^{-q}$, while its Dirichlet character contributes the
  eigenvalue $\chi(p)$.
\end{enumerate}

These statements sit next to, but are different from, several established theories.
Difference-ring algorithms handle very general nested and cyclotomic sums
\cite{Schneider2016,ABS2011}; the point here is an elementary finite construction that
lands directly in proper hypergeometric creative telescoping.  Parameter deformations
are central in the work of Straub--Zudilin on binomial powers and Ap\'ery limits
\cite{StraubZudilin}.  Taylor coefficients also occur in refined Lucas congruences for
Ap\'ery numbers \cite{RowlandYassawiKratt}.  Our focus is the inverse lifting operation
and the character-resolved coefficientwise descent of an arbitrary homogeneous harmonic
weight.

Parameter differentiation as a source of harmonic identities is classical and remains
active; representative references include derivative-operator telescoping
\cite{WeiGong}, all-order hypergeometric parameter expansions \cite{KalmykovWardYost},
and the recent monoidal-alphabet organisation of generalized harmonic sums
\cite{Phadikar}.  We have not found in these sources the Boolean connected lift
\eqref{eq:Rdef}, the resulting arbitrary-weight inverse theorem, or the
character-resolved Lucas theorem in the form proved below.  This is a literature-search
boundary, not a claim that no equivalent formulation can exist.

The last section gives a concrete test.  The recurrence
\begin{equation}\label{eq:Erec-intro}
 (n+1)^2u_{n+1}=(12n^2+12n+4)u_n-32n^2u_{n-1}
\end{equation}
has a binomial first solution and a second-solution limit involving Catalan's constant.
The lifting calculus explains, in one formula, both the product of weight-one letters and
the apparently exceptional weight-two correction in a newly found closed form.
Existing Lucas and Gessel--Lucas results for the sporadic families concern their integral
first solutions \cite{MalikStraub,StraubGessel}.  We have not found in the sources
searched the finite $\chi_{-4}$-harmonic second solution below or a twisted Lucas law for
it; as above, this records the boundary of our search rather than asserting absolute
priority.

\paragraph{Status convention.}
Every theorem and lemma below is proved.  The Catalan recurrence proof is
computer-assisted but exact: its scalar recurrences, right divisions over $\Q(x)$,
leading coefficient polynomials, and finite propagation bounds are reproduced below and
in the accompanying scripts.  Irrationality of Catalan's constant remains open; no
transcendence or irrationality conclusion is asserted.

\section{Harmonic alphabets and parameter products}

For $m,q\ge0$ put $H_0^{(q)}=0$ and, for $q\ge1$,
\[
 H_m^{(q)}=\sum_{j=1}^m j^{-q}.
\]
If $\chi$ is periodic, define $H_m^{(q,\chi)}$ by the same formula with numerator
$\chi(j)$.  We give a letter $H^{(q,\chi)}$ weight $q$; a product has total weight the sum
of its letter weights.

The basic parameter product is
\begin{equation}\label{eq:Xdef}
 X_m(t):=\prod_{j=1}^m(1-t/j)^{-1}=\frac{m!}{(1-t)_m}.
\end{equation}
Its logarithm is the harmonic generating series
\begin{equation}\label{eq:logX}
 \log X_m(t)=\sum_{q\ge1}\frac{H_m^{(q)}}q\,t^q.
\end{equation}
For a residue class $a\bmod M$ we similarly set
\begin{equation}\label{eq:Xres}
 X_{m;a,M}(t):=
 \prod_{\substack{1\le j\le m\\j\equiv a\ (M)}}(1-t/j)^{-1}.
\end{equation}
Then $\partial_t^q\log X_{m;a,M}(0)=(q-1)!$ times the harmonic sum restricted to that
residue class.

\subsection{The connected lift}

Let $[r]=\{1,\ldots,r\}$, write $t_T=\sum_{i\in T}t_i$, and take $X_m(0)=1$.

\begin{definition}[connected lift]\label{def:connected}
For $r\ge1$ define
\begin{equation}\label{eq:Rdef}
 \mathcal R_{m,r}(t_1,\ldots,t_r)
 :=\prod_{T\subseteq[r]}X_m(t_T)^{(-1)^{r-|T|}}.
\end{equation}
\end{definition}

For example,
\[
 \mathcal R_{m,1}(u)=X_m(u),\qquad
 \mathcal R_{m,2}(u,v)=\frac{X_m(u+v)}{X_m(u)X_m(v)}.
\]
The next theorem is the elementary engine of the paper.

\begin{theorem}[single-letter inverse lift]\label{thm:singlelift}
For every $m\ge0$ and $r\ge1$,
\begin{equation}\label{eq:singlelift}
 \left.\partial_{t_1}\cdots\partial_{t_r}
 \mathcal R_{m,r}(t_1,\ldots,t_r)\right|_{t=0}
 =(r-1)!H_m^{(r)}.
\end{equation}
Moreover, every nonconstant monomial of $\log\mathcal R_{m,r}$ contains every one of the
variables $t_1,\ldots,t_r$.
\end{theorem}

\begin{proof}
From \eqref{eq:logX},
\[
 \log\mathcal R_{m,r}
 =\sum_{q\ge1}\frac{H_m^{(q)}}q
   \sum_{T\subseteq[r]}(-1)^{r-|T|}t_T^q.
\]
If a monomial omits $t_i$, its coefficient cancels between $T$ and $T\cup\{i\}$.
Consequently every surviving monomial contains all $r$ variables and has degree at least
$r$.  The coefficient of $t_1\cdots t_r$ comes only from $q=r$; in $t_T^r$ it is zero
unless $T=[r]$, when it is $r!$.  Thus the full mixed derivative of
$\log\mathcal R_{m,r}$ at the origin is $(r-1)!H_m^{(r)}$.  Since the logarithm has no
terms of degree below $r$, the same mixed derivative of its exponential
$\mathcal R_{m,r}$ has the same value.
\end{proof}

This is a finite-difference version of the moment--cumulant relation.  The Boolean signs
are exactly what prevents, for example, $H_m^2$ from contaminating the lift of
$H_m^{(2)}$.

\begin{theorem}[inverse lifting of a harmonic polynomial]\label{thm:universal}
Let $W$ be a polynomial over a characteristic-zero field in finitely many letters
$H_{m_i}^{(q)}$, homogeneous of total weight $w$.  Then $W$ is a finite linear
combination of order-$w$ mixed derivatives at the origin of products and quotients of
terms $X_{m_i}(L(t))$, where $L$ is a linear form with coefficients $0$ or $1$.

The construction is explicit: replace each occurrence of $H_{m}^{(q)}$ by
$\mathcal R_{m,q}/(q-1)!$, give different occurrences disjoint parameter blocks, multiply,
and differentiate once in every parameter.
\end{theorem}

\begin{proof}
It suffices by linearity to treat a monomial
$\prod_{i=1}^s H_{m_i}^{(q_i)}$ with $\sum_iq_i=w$.  Give its $i$th factor a fresh block
of $q_i$ variables and multiply the connected lifts
$\mathcal R_{m_i,q_i}/(q_i-1)!$.  The variables are disjoint, so the order-$w$ derivative
factors.  Theorem~\ref{thm:singlelift} gives the required monomial.
\end{proof}

\begin{corollary}[economical weight two]\label{cor:weighttwo}
Every quadratic harmonic expression in the letters $h_i=H_{m_i}$ and
$q_i=H_{m_i}^{(2)}$ is a linear combination of mixed derivatives of
\begin{equation}\label{eq:weight2blocks}
 X_{m_i}(u)X_{m_j}(v),\qquad
 \frac{X_{m_i}(u+v)}{X_{m_i}(u)X_{m_i}(v)}.
\end{equation}
Indeed the first block gives $h_ih_j$ and the second gives $q_i$.
\end{corollary}

\begin{remark}[polarisation]
If a two-parameter term $F(u,v)$ is expensive to telescope, its mixed derivative can be
recovered from three one-parameter families:
\begin{equation}\label{eq:polarization}
 F_{uv}(0,0)=\frac12\left(\frac{d^2}{dt^2}F(t,t)\big|_{t=0}
 -F_{uu}(0,0)-F_{vv}(0,0)\right).
\end{equation}
This elementary identity was effective in the binomial companion calculations that
motivated the present formalism.
\end{remark}

\section{Cyclotomic letters and finite residue splitting}

Let $M\ge1$.  Every periodic weight $\chi$ of period $M$ decomposes into residue
indicators, so
\begin{equation}\label{eq:resdecomp}
 H_m^{(q,\chi)}=
 \sum_{a=1}^M\chi(a)
 \sum_{\substack{1\le j\le m\\j\equiv a\ (M)}}j^{-q}.
\end{equation}
Applying Theorem~\ref{thm:singlelift} to $X_{m;a,M}$ therefore gives an inverse lift for
every cyclotomic letter.  The construction is over the field generated by the values of
$\chi$.

The relevance to hypergeometric summation is that residue products are Pochhammer ratios.
Write $m=M\ell+s$, $0\le s<M$.  Up to a factor independent of $\ell$,
\begin{equation}\label{eq:respoch}
 X_{M\ell+s;a,M}(t)
 =\frac{(a/M)_{\ell+\epsilon_{a,s}}}
 {((a-t)/M)_{\ell+\epsilon_{a,s}}},
 \qquad \epsilon_{a,s}\in\{0,1\},
\end{equation}
with the evident convention for the representative $a\in\{1,\ldots,M\}$.  Its quotient
under $\ell\mapsto\ell+1$ is rational in $\ell$ and $t$.

\begin{theorem}[cyclotomic inverse lift]\label{thm:cyclolift}
Let $W$ be a polynomial in finitely many cyclotomic harmonic letters with a common period
dividing $M$.  After splitting each affine harmonic argument and each summation index
modulo $M$, $W$ is a finite linear combination of parameter derivatives of products of
Pochhammer ratios whose shift quotients are rational.
\end{theorem}

\begin{proof}
Use \eqref{eq:resdecomp}, Theorem~\ref{thm:universal}, and then
\eqref{eq:respoch}.  Only finitely many residue choices occur.  On each choice, floors and
the corrections $\epsilon_{a,s}$ are fixed affine functions, so all shift quotients are
rational.
\end{proof}

For a real primitive character one may package the residue factors more compactly as
\begin{equation}\label{eq:Ychi}
 Y_m^\chi(t):=\prod_{j=1}^m(1-t/j)^{-\chi(j)},\qquad
 \log Y_m^\chi(t)=\sum_{q\ge1}\frac{H_m^{(q,\chi)}}q t^q.
\end{equation}
When $\chi(j)\in\{0,\pm1\}$ this itself has rational shift quotients after residue
splitting.  For non-real characters the residue-indicator construction is preferable: it
keeps every hypergeometric exponent integral.

\section{Holonomicity is automatic}

We use ``proper hypergeometric'' in the standard Wilf--Zeilberger sense: the quotients
under unit shifts of the discrete variables are rational, with finite factorial support.

\begin{theorem}[cyclotomic harmonic holonomicity]\label{thm:holonomic}
Let $S(n,k)$ be a proper hypergeometric term with finite, linearly bounded support in $k$.
Let $W(n,k)$ be a polynomial in cyclotomic harmonic letters whose arguments are
nonnegative affine integer forms on that support.  Then
\begin{equation}\label{eq:BW}
 B(n)=\sum_kS(n,k)W(n,k)
\end{equation}
is P-recursive over the number field generated by the coefficients and character values.
Moreover, a recurrence and certificate for $B$ can be obtained by finitely many ordinary
proper-hypergeometric creative-telescoping computations followed by linear elimination.
\end{theorem}

\begin{proof}
Choose a common period $M$ and apply Theorem~\ref{thm:cyclolift}.  Split $n,k$ modulo $M$.
Each resulting lift is a proper hypergeometric term over a rational function field in its
parameters.  Zeilberger's theorem supplies a recurrence for its finite definite sum
\cite{WZ1990,PWZ}.  Differentiate the recurrence to the finite order prescribed by $W$
and set the parameters to zero.  A derivative of a parameter-dependent recurrence may
couple the target derivative to lower derivatives, but there are only finitely many of
them.  Closure of holonomic systems under finite elimination gives a scalar recurrence
for the desired derivative.  Finally interlace the finitely many residue subsequences;
P-recursive sequences are closed under finite interlacing.
\end{proof}

\begin{remark}[what the theorem does and does not optimize]
Theorem~\ref{thm:holonomic} guarantees an algorithm, not a minimal recurrence.  Direct
difference-field summation can be dramatically smaller, and parameter families may
produce large nonminimal operators before elimination.  The advantage is universality:
failure of a direct harmonic Gosper ansatz is never a failure of holonomicity.
\end{remark}

\begin{corollary}[bounded lifting complexity]\label{cor:complexity}
If $W$ is a sum of $N$ harmonic monomials of total weight $w$, the construction needs at
most $N$ order-$w$ parameter derivatives.  At weight two, one may instead use exactly one
block of \eqref{eq:weight2blocks} per nonzero quadratic coefficient.  Thus parameter
lifting has complexity controlled by the support of $W$, not by the sizes of its harmonic
arguments.
\end{corollary}

\section{Character-resolved Lucas descent}

Fix a prime $p$.  In this section all $\chi$ are Dirichlet characters (the trivial
character is allowed).  Congruences involving rational harmonic numbers are interpreted in
$\Zp$ after the displayed power of $p$ has cleared the possible $p$-denominators.  Assume
throughout this section that $p$ does not divide the periods of the characters.

\begin{lemma}[one-letter descent]\label{lem:letterdescent}
Let $q\ge1$, let $\chi$ be completely multiplicative, and let $y\ge0$.  Then
\begin{equation}\label{eq:letterdescent}
 p^q H_y^{(q,\chi)}
 =\chi(p)H_{\lfloor y/p\rfloor}^{(q,\chi)}+p^qT,
 \qquad
 T=\sum_{\substack{j\le y\\p\nmid j}}\frac{\chi(j)}{j^q}\in\Zp.
\end{equation}
\end{lemma}

\begin{proof}
Separate the summands with $p\mid j$, write $j=p\ell$, and use
$\chi(p\ell)=\chi(p)\chi(\ell)$.  The complementary denominators are $p$-units.
\end{proof}

A character signature is a product $\eta=\prod_i\chi_i$, with repetitions allowed.  A
monomial $\prod_iH_{x_i}^{(q_i,\chi_i)}$ has weight $w=\sum_iq_i$ and signature $\eta$.

\begin{lemma}[monomial descent]\label{lem:monomial}
Suppose $0\le y_i<p^2$ for every $i$.  Then
\begin{equation}\label{eq:monomial}
 p^w\prod_iH_{y_i}^{(q_i,\chi_i)}
 \equiv \eta(p)\prod_iH_{\lfloor y_i/p\rfloor}^{(q_i,\chi_i)}\pmod p.
\end{equation}
Both sides belong to $\Zp$.
\end{lemma}

\begin{proof}
Multiply the identities \eqref{eq:letterdescent}.  Because
$\lfloor y_i/p\rfloor<p$, every descended denominator is a $p$-unit.  Every term other
than the product of the leading terms contains a positive power of $p$.
\end{proof}

We now state a strong, clean digit theorem.  It is designed so that no informal
``surviving rectangle'' argument is needed.

\begin{theorem}[character-resolved harmonic Lucas theorem]\label{thm:lucas}
Let $S(n,k)\in\Z$, extended by zero outside its finite support, and put
$A(n)=\sum_kS(n,k)$.  Let
\[
 W=\sum_\eta W_\eta
\]
be a finite sum in which each $W_\eta(n,k)$ is a homogeneous weight-$w$ polynomial in
cyclotomic harmonic letters of signature $\eta$.  Put
$B_\eta(n)=\sum_kS(n,k)W_\eta(n,k)$.

Let $n=ap+r$ with $0\le a,r<p$.  Assume:
\begin{enumerate}[label=\textup{(L\arabic*)},leftmargin=2.4em]
\item for every $b\ge0$ and $0\le s<p$,
\[
 S(ap+r,bp+s)\equiv S(a,b)S(r,s)\pmod p;
\]
\item every harmonic argument $x$ occurring in $W$ satisfies
$0\le x(ap+r,bp+s)<p^2$ on the support and
\[
 \left\lfloor\frac{x(ap+r,bp+s)}p\right\rfloor=x(a,b);
\]
\item the coefficients of $W$ belong to $\Zp$.
\end{enumerate}
Then
\begin{equation}\label{eq:lucasmain}
 p^wB_\eta(ap+r)\equiv \eta(p)B_\eta(a)A(r)\pmod p
\end{equation}
for every signature $\eta$, and hence
\begin{equation}\label{eq:lucassum}
 p^wB(ap+r)\equiv A(r)\sum_\eta\eta(p)B_\eta(a)\pmod p.
\end{equation}
\end{theorem}

\begin{proof}
Write each $k$ uniquely as $bp+s$, $0\le s<p$, and enlarge to a full finite block using
the zero extension of $S$.  Lemma~\ref{lem:monomial}, linearity, and (L2) give
\[
 p^wW_\eta(ap+r,bp+s)\equiv\eta(p)W_\eta(a,b)\pmod p.
\]
This expression is $p$-integral on every cell.  Therefore (L1) gives
\begin{align*}
 p^wB_\eta(ap+r)
 &\equiv\eta(p)\sum_{b,s}S(a,b)S(r,s)W_\eta(a,b)\pmod p\\
 &=\eta(p)B_\eta(a)A(r),
\end{align*}
which proves \eqref{eq:lucasmain}; summing over signatures gives
\eqref{eq:lucassum}.
\end{proof}

\begin{corollary}[Frobenius eigenspaces]\label{cor:eigenspaces}
If all monomials have one signature $\eta$, the companion has the scalar digit law
\eqref{eq:lucasmain}.  If several signatures occur, their components do not mix: the
digit map acts diagonally with eigenvalues $\eta(p)$.  Thus character-homogeneity is not a
technical convenience but the condition that the harmonic jet lie in one Frobenius
eigenspace.
\end{corollary}

\begin{corollary}[coefficientwise jet descent]\label{cor:jet}
Collect all homogeneous harmonic polynomials of weight $w$ satisfying (L2) into a formal
weight-$w$ jet $\mathcal J_w(n)$.  Then the digit map on $\mathcal J_w$ is the tensor
product of multiplication by $A(r)$ with the diagonal character action
$\eta\mapsto\eta(p)$, together with the Tate factor $p^{-w}$.
\end{corollary}

\begin{remark}[scope]
The restriction $a<p$ ensures that descended harmonic numbers are $p$-integral.  Iterating
\eqref{eq:lucasmain} without rechecking this point is not justified.  Multi-digit laws
require either stronger denominator control or a divided-power formulation.  The theorem
is deliberately single-digit.
\end{remark}

\section{A compact lift for the Catalan companion}

Let $\chi_{-4}$ be the real character modulo $4$:
\[
 \chi_{-4}(j)=0\ (2\mid j),\qquad
 \chi_{-4}(j)=(-1)^{(j-1)/2}\ (j\text{ odd}),
\]
and abbreviate $K_m^{(q)}=H_m^{(q,\chi_{-4})}$.  Define
\begin{equation}\label{eq:Eshell}
 S_E(n,k)=\binom nk\binom{2k}k\binom{2n-2k}{n-k},\qquad
 A_E(n)=\sum_{k=0}^nS_E(n,k).
\end{equation}
The sequence $A_E=1,4,20,112,676,\ldots$ satisfies \eqref{eq:Erec-intro}.  The second
solution normalized by $B_E(0)=0$, $B_E(1)=1$ has Ap\'ery limit $G/2$, where
$G=L(2,\chi_{-4})$ is Catalan's constant; this value belongs to the established sporadic
Ap\'ery-limit tables \cite{Zagier,ChamStraub,MalikStraub}.
This $A_E$ is the scaled Catalan--Larcombe--French sequence.  A useful second shell is
\begin{equation}\label{eq:parityshell}
 A_E(n)=\sum_{0\le k\le n/2}
 \binom{2k}{k}^{\!2}\binom n{2k}4^{n-2k},
\end{equation}
and its first-solution congruences have been studied extensively
\cite{JiSunCLF,StraubGessel}.  The arithmetic results below instead concern its rational
second solution.

\begin{theorem}[Catalan harmonic companion]\label{thm:Catalan}
For every $n\ge0$, the normalized second solution of \eqref{eq:Erec-intro} is
\begin{equation}\label{eq:CatalanB}
 B_E(n)=\sum_{k=0}^nS_E(n,k)\left\{
 \frac12K_{2k}^{(2)}+
 \left(\frac34H_k-\frac12H_{2k}\right)
 \left(K_{2k}^{(1)}-K_{2n-2k}^{(1)}\right)\right\}.
\end{equation}
The first values are
\[
 0,\ 1,\ 7,\ \frac{404}{9},\ \frac{2603}{9},\
 \frac{428284}{225},\ \frac{2884316}{225}.
\]
These values satisfy \eqref{eq:Erec-intro} for every $n\ge1$.  Moreover,
\begin{equation}\label{eq:Catalanlimit}
 \lim_{n\to\infty}\frac{B_E(n)}{A_E(n)}=\frac G2.
\end{equation}
\end{theorem}

The formula was originally found by exact linear algebra in a weight-two cyclotomic
alphabet.  The following proposition explains why all three pieces of its weight occur
together.  The recurrence proof is completed after the generating-function
factorisation below.

\begin{proposition}[two-parameter Catalan lift]\label{prop:Catalanlift}
Let $X$ be \eqref{eq:Xdef}, let $Y_m(t)=Y_m^{\chi_{-4}}(t)$ be \eqref{eq:Ychi}, and put
\begin{align}\label{eq:FE}
 \mathcal F_E(n,k;u,v):={}&S_E(n,k)
 \frac{X_k(3u/4)}{X_{2k}(u/2)}
 \frac{Y_{2k}(v)}{Y_{2n-2k}(v)}\\[-2mm]
 &\hspace{8mm}\times
 \frac{Y_{2k}(u+v/2)}{Y_{2k}(u)Y_{2k}(v/2)}.\notag
\end{align}
Then the summand in \eqref{eq:CatalanB} is exactly
\begin{equation}\label{eq:mixedFE}
 \left.\partial_u\partial_v\mathcal F_E(n,k;u,v)\right|_{u=v=0}.
\end{equation}
\end{proposition}

\begin{proof}
All parameter factors equal $1$ at the origin.  From \eqref{eq:logX} and
\eqref{eq:Ychi}, the first logarithmic derivatives of \eqref{eq:FE} are
\[
 \partial_u\log\mathcal F_E\big|_0=\frac34H_k-\frac12H_{2k},\qquad
 \partial_v\log\mathcal F_E\big|_0=K_{2k}^{(1)}-K_{2n-2k}^{(1)}.
\]
Only the last quotient in \eqref{eq:FE} has a mixed logarithmic derivative.  It is the
connected lift $Y_{2k}(u+v/2)/(Y_{2k}(u)Y_{2k}(v/2))$, so
\[
 \partial_u\partial_v\log\mathcal F_E\big|_0=\frac12K_{2k}^{(2)}.
\]
Finally $F_{uv}/F=(\log F)_u(\log F)_v+(\log F)_{uv}$, which is precisely the weight in
\eqref{eq:CatalanB}.
\end{proof}

The lift is not merely formal: it reduces to proper hypergeometric terms.

\begin{lemma}[parity reduction]\label{lem:parity}
Put $P_q(t)=Y_{2q}(t)$.  For $j\ge0$,
\begin{align}
 P_{2j}(t)
 &=\frac{(1/4)_j}{((1-t)/4)_j}
   \frac{((3-t)/4)_j}{(3/4)_j},\label{eq:P2j}\\
 P_{2j+1}(t)
 &=P_{2j}(t)\frac{4j+1}{4j+1-t}.\label{eq:P2j1}
\end{align}
\end{lemma}

\begin{proof}
In $Y_{4j}$ the positive character values occur at $1,5,\ldots,4j-3$ and the negative
ones at $3,7,\ldots,4j-1$.  Grouping these two products and extracting powers of $4$
gives \eqref{eq:P2j}.  Passing from $4j$ to $4j+2$ adds only the positive factor
$(1-t/(4j+1))^{-1}$, giving \eqref{eq:P2j1}.
\end{proof}

\begin{corollary}[certifiability of the Catalan formula]\label{cor:certCatalan}
After writing $n=2m+\epsilon$ and $k=2j+\delta$, the summand
$\mathcal F_E(n,k;u,v)$ is proper hypergeometric in $m,j$ on each of the four parity
cells.  Consequently the right side of \eqref{eq:CatalanB} is P-recursive and admits a
finite creative-telescoping certificate.
\end{corollary}

\begin{proof}
The $X$-factors are factorial ratios.  Lemma~\ref{lem:parity} turns every $Y$-factor into
a Pochhammer ratio after the parity choice.  The remaining factor $S_E$ is proper
hypergeometric.  Apply Theorem~\ref{thm:holonomic}.
\end{proof}

There is also a useful convolution structure.  With $\ell=n-k$, \eqref{eq:FE} factors as
\[
 \binom nk\,P_k(u,v)Q_\ell(v),
\]
where
\begin{align*}
 P_k(u,v)&=\binom{2k}k\frac{X_k(3u/4)}{X_{2k}(u/2)}
 \frac{Y_{2k}(v)Y_{2k}(u+v/2)}{Y_{2k}(u)Y_{2k}(v/2)},\\
 Q_\ell(v)&=\binom{2\ell}\ell\,Y_{2\ell}(v)^{-1}.
\end{align*}
Thus the exponential generating function of the lifted sum is the product
\begin{equation}\label{eq:EGFfactor}
 \left(\sum_{k\ge0}P_k(u,v)\frac{x^k}{k!}\right)
 \left(\sum_{\ell\ge0}Q_\ell(v)\frac{x^\ell}{\ell!}\right).
\end{equation}
The determinant hidden in this factorisation leads to a substantially smaller proof than
the raw parity-split telescoper.  Put
\begin{equation}\label{eq:cadef}
 c_k=\binom{2k}{k},\qquad
 \alpha_k=\frac34H_k-\frac12H_{2k},\qquad
 a_k=c_k\alpha_k,\qquad d_k=c_kK_{2k}^{(1)},
\end{equation}
and
\begin{equation}\label{eq:edef}
 e_k=c_k\left(\frac12K_{2k}^{(2)}+\alpha_kK_{2k}^{(1)}\right).
\end{equation}
If capital letters denote exponential generating functions, binomial convolution gives
\begin{equation}\label{eq:Bdet}
 B(x):=\sum_{n\ge0}B_E(n)\frac{x^n}{n!}=E(x)C(x)-A(x)D(x).
\end{equation}
Thus the companion is a $2\times2$ determinant of four harmonic jets.

The elementary increments behind the certificate are
\begin{align}
 \alpha_{k+1}-\alpha_k&=\frac{k}{2(k+1)(2k+1)},\notag\\
 K_{2k+2}^{(1)}-K_{2k}^{(1)}&=\frac{(-1)^k}{2k+1},&
 K_{2k+2}^{(2)}-K_{2k}^{(2)}&=\frac{(-1)^k}{(2k+1)^2}.
 \label{eq:increments}
\end{align}
With $p_k=(-1)^kc_k$ and $s_k=(-1)^ka_k$, direct substitution gives the coupled system
\begin{align}
 (k+1)c_{k+1}&=2(2k+1)c_k,\notag\\
 (k+1)^2a_{k+1}&=2(k+1)(2k+1)a_k+kc_k,\notag\\
 (k+1)d_{k+1}&=2(2k+1)d_k+2p_k,\notag\\
 (k+1)p_{k+1}&=-2(2k+1)p_k,\label{eq:coupledsystem}\\
 (k+1)^2s_{k+1}&=-2(k+1)(2k+1)s_k-kp_k,\notag\\
 (k+1)^2e_{k+1}&=2(k+1)(2k+1)e_k+p_k+2(k+1)s_k+kd_k.\notag
\end{align}
In particular, the scalar input recurrences below are consequences of first-order
rational identities, not interpolated recurrences.

Let $\theta=xD_x$ and define
\begin{equation}\label{eq:LaperyEGF}
 \mathcal L=(\theta+1)^2D_x-(12\theta^2+12\theta+4)+32x(\theta+1).
\end{equation}
For an exponential generating function $F(x)=\sum f_nx^n/n!$, the coefficient of
$x^n/n!$ in $\mathcal LF$ is
\[
 (n+1)^2f_{n+1}-(12n^2+12n+4)f_n+32n^2f_{n-1}.
\]
Set $Y_1=\mathcal L(EC)$ and $Y_2=\mathcal L(AD)$.

\begin{proposition}[exact common-factor certificate]\label{prop:OreCatalan}
The two residual series satisfy
\begin{equation}\label{eq:Ydifference}
 Y_1-Y_2=1.
\end{equation}
Exact holonomic closure applied to \eqref{eq:coupledsystem} produces annihilators
$\mathcal H_1,\mathcal H_2$ of orders $44,29$.  After clearing scalar denominators,
both have the same right factor $\mathcal C\in\Q[x]\langle D_x\rangle$, of order $10$
and degree $29$:
\begin{equation}\label{eq:rightfactors}
 \mathcal H_1=\mathcal U_1\mathcal C,\qquad
 \mathcal H_2=\mathcal U_2\mathcal C,\qquad
 \ord(\mathcal U_1)=34,\quad\ord(\mathcal U_2)=19.
\end{equation}
Moreover $\mathcal C(1)=0$.
\end{proposition}

\begin{proof}
All operations are exact over $\Q(x)$.  Elimination from
\eqref{eq:coupledsystem}, exponential normalisation, product closure, and association by
\eqref{eq:LaperyEGF} give $\mathcal H_1$ and $\mathcal H_2$.  An independently fitted
order-$11$, degree-$17$ operator is used only to locate $\mathcal C$; the identities in
\eqref{eq:rightfactors} are then established by exact zero-remainder right division, so
the fit has no logical role in the certificate.

We record the finite uniqueness check.  For a polynomial differential operator
$T=\sum t_{j,i}x^iD_x^j$, the coefficient of $f_{n+r}$ in $T(\sum f_nx^n)$, where
$r=\max(j-i)$, is
\begin{equation}\label{eq:forwardcoeff}
 P_T(n)=\sum_{j-i=r}t_{j,i}(n-i+1)\cdots(n-i+j).
\end{equation}
For $\mathcal U_1,\mathcal U_2,\mathcal C$, the largest forward shift is $4$ and, up to
nonzero rational constants, the three forward coefficients are
\begin{align*}
 &(z+1)(z+2)^2(z+3)^2(z+4)^3(z+5)^3\prod_{j=6}^{9}(z+j)^2\\[-1mm]
 &\quad\times
 \left(z^5+\frac{61}{2}z^4+370z^3+\frac{4461}{2}z^2
             +\frac{26709}{4}z+\frac{63465}{8}\right)^3,\\
 &(z+1)(z+2)^2(z+3)^2(z+4)^3(z+5)^3\prod_{j=6}^{9}(z+j)^2,\\
 &(z+1)(z+2)^2(z+3)^2(z+4)^4(z+5).
\end{align*}
They are nonzero at every nonnegative integer.  Their complete forward coefficients are
present from $z=30,15,6$, respectively.  Exact expansion gives the first $34$
coefficients of $\mathcal CY_1$ and the first $19$ coefficients of $\mathcal CY_2$ equal
to zero.  The first two recurrences therefore imply
$\mathcal CY_1=\mathcal CY_2=0$.  Since $\mathcal C(1)=0$, the series $Y_1-Y_2-1$ is a
$\mathcal C$-solution; its first $10$ coefficients vanish, and the third recurrence gives
\eqref{eq:Ydifference}.

The reconstructed operator $\mathcal C$ has SHA-256 fingerprint
\[
 \texttt{5f255471c10932d5a24a7cf4f429729bda8d9f73740aa3a1c37e261e70312fba}.
\]
The accompanying exact-arithmetic scripts reconstruct it and check every assertion.
\end{proof}

\begin{proof}[Proof of Theorem~\ref{thm:Catalan}, recurrence part]
By \eqref{eq:Bdet}, Proposition~\ref{prop:OreCatalan} says $\mathcal LB=1$.
Coefficient comparison gives \eqref{eq:Erec-intro} for $n\ge1$.  The finite sum gives
$B_E(0)=0$ and $B_E(1)=1$, hence it is the normalized second solution.
\end{proof}

The harmonic representation also evaluates its own Ap\'ery limit, without appealing to
the sporadic tables.  This is one of the principal benefits of having a finite companion
formula.

\begin{proposition}[direct Catalan limit]\label{prop:directlimit}
Formula \eqref{eq:CatalanB} implies \eqref{eq:Catalanlimit} by elementary estimates.
\end{proposition}

\begin{proof}
Write $D_m=K_{2m}^{(1)}$, $Q_m=K_{2m}^{(2)}$, and put
$\mu_n(k)=S_E(n,k)/A_E(n)$.  Since $S_E(n,k)=S_E(n,n-k)$, symmetrising
\eqref{eq:CatalanB} gives, with $\ell=n-k$,
\begin{equation}\label{eq:symaverage}
 \frac{B_E(n)}{A_E(n)}=
 \sum_{k=0}^n\mu_n(k)\left\{
 \frac14(Q_k+Q_\ell)
 +\frac12(\alpha_k-\alpha_\ell)(D_k-D_\ell)\right\}.
\end{equation}

We first control the shell measure.  The elementary bounds
\[
 \binom{2m}{m}\le4^m,\qquad
 \binom{2m}{m}\ge\frac{4^m}{2m+1},\qquad
 \max_k\binom nk\ge\frac{2^n}{n+1}
\]
give, by taking $k=\lfloor n/2\rfloor$,
\begin{equation}\label{eq:Alower}
 A_E(n)\ge\frac{8^n}{(n+1)^3}.
\end{equation}
On the other hand, $S_E(n,k)\le4^n\binom nk$.  The exponential-moment estimate
\[
 \sum_{k\le n/4}\binom nk
 \le 3^{n/4}\sum_{k=0}^n\binom nk3^{-k}
 =\frac{4^n}{3^{3n/4}}
\]
and symmetry therefore show that
\begin{equation}\label{eq:boundarymass}
 \sum_{k<n/4\ \mathrm{or}\ k>3n/4}\mu_n(k)
 \le2(n+1)^3\left(\frac{2}{3^{3/4}}\right)^n.
\end{equation}

On the bulk $n/4\le k,\ell\le3n/4$, the alternating-series estimate gives
\begin{equation}\label{eq:alternatingbounds}
 |Q_m-G|\le(2m+1)^{-2},\qquad
 |D_k-D_\ell|\le(2\min(k,\ell)+1)^{-1}.
\end{equation}
Moreover \eqref{eq:increments} implies
\[
 0\le\alpha_{m+1}-\alpha_m\le\frac1{4(m+1)},
\]
so $|\alpha_k-\alpha_\ell|$ is bounded uniformly on the bulk.  Thus the expression in
braces in \eqref{eq:symaverage} is $G/2+O(1/n)$ there.  Globally,
$|D_m|,|Q_m|\le1$ and $0\le\alpha_m\le H_m/4$, so the same expression is
$O(\log(n+1))$.  Equation \eqref{eq:boundarymass} makes the boundary contribution tend
to zero, proving \eqref{eq:Catalanlimit}.
\end{proof}

\begin{proof}[Completion of the proof of Theorem~\ref{thm:Catalan}]
The recurrence and initial values were proved above, and
Proposition~\ref{prop:directlimit} proves the claimed limit.
\end{proof}

\begin{corollary}[positive error formula]\label{cor:Catalanerror}
The approximants $B_E(n)/A_E(n)$ increase strictly to $G/2$, and
\begin{align}
 \frac G2-\frac{B_E(n)}{A_E(n)}
 &=\sum_{m=n+1}^{\infty}
   \frac{32^{m-1}}{m^2A_E(m)A_E(m-1)},\label{eq:errorseries}\\
 G&=2\sum_{m=1}^{\infty}
   \frac{32^{m-1}}{m^2A_E(m)A_E(m-1)}.\label{eq:Gseries}
\end{align}
\end{corollary}

\begin{proof}
For two solutions $A,B$ of \eqref{eq:Erec-intro}, their Casoratian
$W_n=A_nB_{n-1}-A_{n-1}B_n$ satisfies
\[
 W_{n+1}=\frac{32n^2}{(n+1)^2}W_n.
\]
The initial values give $W_1=-1$, hence $W_n=-32^{n-1}/n^2$.  Dividing by
$A_E(n)A_E(n-1)>0$ yields
\[
 \frac{B_E(n)}{A_E(n)}-\frac{B_E(n-1)}{A_E(n-1)}
 =\frac{32^{n-1}}{n^2A_E(n)A_E(n-1)}>0.
\]
Sum this identity and use Proposition~\ref{prop:directlimit}; taking $n=0$ gives
\eqref{eq:Gseries}.
\end{proof}

The error formula does not expose the true denominator: its summands contain quotients
of first solutions.  The harmonic formula gives a first bound.

\begin{corollary}[termwise denominator bound]\label{cor:Catalanden}
Put $L_0=1$ and $L_m=\operatorname{lcm}(1,2,\ldots,m)$ for $m\ge1$.  Then
\begin{equation}\label{eq:denbound}
 4L_{2n}^2B_E(n)\in\Z.
\end{equation}
\end{corollary}

\begin{proof}
Every first-order harmonic letter in \eqref{eq:CatalanB} has denominator dividing
$L_{2n}$, and every second-order letter has denominator dividing $L_{2n}^2$.
The factor $4$ clears the rational coefficients in $\alpha_k$ and the leading $1/2$;
the shell is integral.
\end{proof}

The factor $L_{2n}^2$ is not sharp.  A binomial transform at the endpoint $4$ of the
two characteristic roots exposes all the cancellation that is invisible termwise in
\eqref{eq:CatalanB}.

\begin{theorem}[sharp global denominator]\label{thm:sharpden}
For every $n\ge0$,
\begin{equation}\label{eq:sharpden}
 L_n^2B_E(n)\in\Z.
\end{equation}
More precisely, define the endpoint transform
\begin{equation}\label{eq:Ttransform}
 T_n=\sum_{k=0}^n\binom nk(-4)^{n-k}B_E(k).
\end{equation}
Then $T_0=0$ and, with $T_{-1}=0$,
\begin{equation}\label{eq:Trec}
 (n+1)^2T_{n+1}-16n^2T_{n-1}=(-4)^n\qquad(n\ge0).
\end{equation}
For $n\ge1$ this has the signed sum-of-squares solution
\begin{equation}\label{eq:Tsquares}
 (-1)^{n-1}T_n=
 \sum_{\substack{0\le j<n\\n-j\ {
m odd}}}R_{n,j}^2,
 \qquad
 R_{n,j}=2^{n-1}
 \frac{(j+2)(j+4)\cdots(n-1)}{(j+1)(j+3)\cdots n},
\end{equation}
where an empty numerator is $1$.  Every $L_nR_{n,j}$ is an integer.
\end{theorem}

\begin{proof}
Let $B(z)=\sum_{n\ge0}B_E(n)z^n$ and $\theta=zD_z$.  Including the exceptional
initial residual $B_E(1)-4B_E(0)=1$, the recurrence is equivalent to
\begin{equation}\label{eq:BOGF}
 \bigl\{\theta^2-z(12\theta^2+12\theta+4)
       +32z^2(\theta+1)^2\bigr\}B(z)=z.
\end{equation}
The ordinary generating function of \eqref{eq:Ttransform} is
\begin{equation}\label{eq:TOGF}
 T(z)=\frac1{1+4z}B\!\left(\frac z{1+4z}\right).
\end{equation}
A chain-rule substitution in \eqref{eq:BOGF} gives
\[
 \{\theta^2-16z^2(\theta+1)^2\}T(z)=\frac z{1+4z}.
\]
Coefficient comparison is exactly \eqref{eq:Trec}.

Iterating \eqref{eq:Trec} separately on the even and odd indices shows that the
contribution of its forcing term $(-4)^j$ to $T_n$, for $n-j$ odd, is
\[
 \frac{(-4)^j}{(j+1)^2}
 \prod_{\substack{j+2\le r<n\\r\equiv j\ (2)}}
 \frac{16r^2}{(r+1)^2}=(-1)^jR_{n,j}^2.
\]
All admissible $j$ have parity $n-1$, proving \eqref{eq:Tsquares}.

It remains to prove the asserted denominator of $R_{n,j}$.  Fix an odd prime $p$.
For every $a\ge1$, the multiples of $p^a$ in the interval $[j+1,n]$ are $p^a$
times a consecutive interval of integers.  Because $p^a$ is odd, multiplication by it
preserves parity.  Hence the number of those multiples in the denominator progression
of $R_{n,j}$ exceeds the number in its numerator progression by at most one.  Summing
over $a$ gives
\[
 -v_p(R_{n,j})\le \#\{a:p^a\le n\}=v_p(L_n)
\]
whenever the left side is positive.  At $p=2$, the denominator progression is either
odd or has $2$-adic valuation at most $v_2(n!)\le n-1$, which is absorbed by the
factor $2^{n-1}$.  Thus $L_nR_{n,j}\in\mathbb Z$.

Equation \eqref{eq:Tsquares} now gives $L_n^2T_n\in\mathbb Z$.  Binomial inversion in
\eqref{eq:Ttransform} reads
\[
 B_E(n)=\sum_{k=0}^n\binom nk4^{n-k}T_k.
\]
Since $L_k\mid L_n$ for $k\le n$, every summand becomes integral after multiplication
by $L_n^2$, proving \eqref{eq:sharpden}.
\end{proof}

\begin{theorem}[general endpoint-denominator classification]
\label{thm:generalendpoint}
Let $b\in\mathbb Z$, and let $B^{(b)}_n$ be the unique rational sequence with
$B^{(b)}_0=0$, $B^{(b)}_1=1$, and
\begin{equation}\label{eq:brec}
 (n+1)^2B^{(b)}_{n+1}
 =b(3n^2+3n+1)B^{(b)}_n-2b^2n^2B^{(b)}_{n-1}
 \qquad(n\ge1).
\end{equation}
Define
\begin{equation}\label{eq:btransform}
 T^{(b)}_n=\sum_{k=0}^n\binom nk(-b)^{n-k}B^{(b)}_k.
\end{equation}
Then
\begin{equation}\label{eq:bTrec}
 (n+1)^2T^{(b)}_{n+1}-b^2n^2T^{(b)}_{n-1}=(-b)^n
 \qquad(n\ge0),
\end{equation}
and, for $n\ge1$,
\begin{equation}\label{eq:bTsquares}
 (-1)^{n-1}T^{(b)}_n
 =b^{n-1}\!
 \sum_{\substack{0\le j<n\\n-j\ {\rm odd}}}
 \left\{
 \frac{(j+2)(j+4)\cdots(n-1)}
      {(j+1)(j+3)\cdots n}
 \right\}^{\!2}.
\end{equation}
Moreover, the following are equivalent:
\begin{equation}\label{eq:bclassification}
 4\mid b
 \qquad\Longleftrightarrow\qquad
 L_n^2B^{(b)}_n\in\mathbb Z\quad\hbox{for every }n\ge0.
\end{equation}
\end{theorem}

\begin{proof}
Let $B^{(b)}(z)=\sum_{n\ge0}B^{(b)}_nz^n$.  Including the initial residual
$B^{(b)}_1-bB^{(b)}_0=1$, recurrence \eqref{eq:brec} is equivalent to
\begin{equation}\label{eq:bOGF}
 \{\theta^2-bz(3\theta^2+3\theta+1)
       +2b^2z^2(\theta+1)^2\}B^{(b)}(z)=z.
\end{equation}
The binomial transform has generating function
\[
 T^{(b)}(z)=\frac1{1+bz}
 B^{(b)}\!\left(\frac z{1+bz}\right).
\]
Substitution in \eqref{eq:bOGF} cancels both the middle term and every term of order
$z^3$, leaving
\[
 \{\theta^2-b^2z^2(\theta+1)^2\}T^{(b)}(z)=\frac z{1+bz}.
\]
Coefficient comparison proves \eqref{eq:bTrec}.

As in the proof of Theorem~\ref{thm:sharpden}, iterate the even and odd recurrences.
The contribution of the forcing term with index $j$ to $T^{(b)}_n$ is
\[
 \frac{(-b)^j}{(j+1)^2}
 \prod_{\substack{j+2\le r<n\\r\equiv j\ (2)}}
 \frac{b^2r^2}{(r+1)^2}
 =(-1)^jb^{n-1}
 \left\{
 \frac{(j+2)(j+4)\cdots(n-1)}
      {(j+1)(j+3)\cdots n}
 \right\}^{\!2}.
\]
All admissible $j$ have parity $n-1$, proving \eqref{eq:bTsquares}.

Suppose first that $b=4d$.  Write the rational number in braces in
\eqref{eq:bTsquares} as $P_{n,j}$.  The denominator argument in
Theorem~\ref{thm:sharpden} proved the stronger statement
\[
 L_n\,2^{n-1}P_{n,j}\in\mathbb Z.
\]
Consequently
\[
 L_n^2b^{n-1}P_{n,j}^2
 =d^{\,n-1}\bigl(L_n2^{n-1}P_{n,j}\bigr)^2\in\mathbb Z.
\]
Equation \eqref{eq:bTsquares} therefore gives
$L_n^2T^{(b)}_n\in\mathbb Z$.  Binomial inversion reads
\[
 B^{(b)}_n=\sum_{k=0}^n\binom nkb^{n-k}T^{(b)}_k.
\]
Since $L_k\mid L_n$, multiplying by $L_n^2$ makes every summand integral.  This proves
the forward implication in \eqref{eq:bclassification}.

For the converse, direct propagation of \eqref{eq:brec} gives
\begin{equation}\label{eq:btestvalues}
 B^{(b)}_4=\frac{2603b^3}{576},\qquad
 B^{(b)}_8=\frac{6802537507b^7}{180633600}.
\end{equation}
Since $L_4=12$ and $2603$ is odd, integrality at $n=4$ says
$2603b^3/4\in\mathbb Z$, forcing $b$ to be even.  If $b=2d$ with $d$ odd, then
$L_8=840$ and the second value in \eqref{eq:btestvalues} gives
\[
 L_8^2B^{(b)}_8=\frac{6802537507b^7}{256}
 =\frac{6802537507d^7}{2},
\]
which is not integral.  Thus integrality at every index forces $4\mid b$, completing
the reverse implication.
\end{proof}

\begin{remark}[why this family decouples]\label{rem:endpointlocus}
For the general quadratic recurrence and a nonzero shift $h$,
\[
 (n+1)^2u_{n+1}
 =\{\alpha n(n+1)+\beta\}u_n-\gamma n^2u_{n-1},
\]
the binomial shift conjugates its ordinary generating-function operator to
\begin{align*}
 \theta^2
 &-z\{(\alpha-3h)(\theta^2+\theta)+(\beta-h)\}\\
 &+z^2(\gamma-2\alpha h+3h^2)(\theta+1)^2
 +z^3h(\gamma-\alpha h+h^2)(\theta+1)(\theta+2).
\end{align*}
The coefficient at $z$ vanishes only when $\alpha=3h$ and $\beta=h$.  Under these
conditions the order-three coefficient vanishes exactly when $\gamma=2h^2$.  Thus,
apart from the trivial zero-shift
degeneracy, \eqref{eq:brec} is exactly the endpoint-decoupling locus, not an arbitrary
parameter insertion.  It is also a dilation:
$B^{(b)}_n=b^{n-1}B^{(1)}_n$ for $n\ge1$.  The arithmetic content of
Theorem~\ref{thm:generalendpoint} is that precisely the dilations $4\mid b$ supply the
full $2$-adic factor needed by the parity quotient while the odd-prime deficit is
absorbed by $L_n$.
\end{remark}

\begin{remark}[where the cancellation occurs]\label{rem:sharpden}
The summands of the harmonic formula \eqref{eq:CatalanB} can retain primes between
$n$ and $2n$, so Theorem~\ref{thm:sharpden} is genuinely global in that shell.  The
endpoint transform moves the cancellation into the elementary parity interlacing of
\eqref{eq:Tsquares}; there it is termwise after multiplication by $L_n^2$.
\end{remark}

The same transform retains the Ap\'ery limit in an unexpected way: $G$ becomes the
finite part of a logarithmically divergent reciprocal-binomial sum.  Reciprocal-binomial
accelerations of $G$ have a substantial history \cite{BradleyCatalan}; the point here is
that this particular zero-balanced finite part is forced directly by the endpoint
transform of the companion.

\begin{corollary}[endpoint finite part and a series for $G$]\label{cor:Gfinitepart}
One has the absolutely convergent identity
\begin{equation}\label{eq:Grecipseries}
 \sum_{j=1}^{\infty}\left\{
  \frac{16^{j-1}}{j^2\binom{2j}{j}^2}-\frac{\pi}{16j}\right\}
 =\frac{\pi}{4}\log2-\frac G2.
\end{equation}
Equivalently,
\begin{equation}\label{eq:Grecipseries2}
 G=\frac{\pi}{2}\log2+
 \sum_{j=1}^{\infty}\left\{
  \frac{\pi}{8j}-\frac{2\cdot16^{j-1}}{j^2\binom{2j}{j}^2}\right\}.
\end{equation}
In particular, the even part of the rational endpoint transform satisfies
\begin{equation}\label{eq:Tevenfinitepart}
 \lim_{m\to\infty}\left\{
  \frac{T_{2m}}{\binom{2m}{m}^2}+\frac{\pi}{16}H_m
  +\frac{\pi}{4}\log2\right\}=\frac G2.
\end{equation}
\end{corollary}

\begin{proof}
Write
\[
 t_j=\frac{16^{j-1}}{j^2\binom{2j}{j}^2},\qquad
 F(x)=\sum_{j\ge1}t_jx^j
 =\frac{x}{4}\,{}_3F_2\!\left(\begin{matrix}1,1,1\\3/2,3/2\end{matrix};x\right).
\]
Euler's integral for one numerator--denominator pair, followed by
\[
 {}_2F_1(1,1;3/2;y)=\frac{\arcsin(\sqrt y)}{\sqrt{y(1-y)}}
\]
and the substitution $t=\sin^2\theta$, gives, for $0<x<1$,
\begin{equation}\label{eq:Hintegral}
 {}_3F_2\!\left(\begin{matrix}1,1,1\\3/2,3/2\end{matrix};x\right)
 =\frac1{\sqrt x}\int_0^{\pi/2}
 \frac{\arcsin(\sqrt x\sin\theta)}
      {\sqrt{1-x\sin^2\theta}}\,d\theta.
\end{equation}
Let $K(x)=\int_0^{\pi/2}(1-x\sin^2\theta)^{-1/2}\,d\theta$.  Subtracting
$\pi/2$ in the numerator of \eqref{eq:Hintegral} decomposes its right side as
\[
 \frac1{\sqrt x}\left\{\frac\pi2K(x)+J(x)\right\},
 \quad
 J(x)=\int_0^{\pi/2}
 \frac{\arcsin(\sqrt x\sin\theta)-\pi/2}
      {\sqrt{1-x\sin^2\theta}}\,d\theta.
\]
The standard endpoint estimate is
$K(x)=\log4-\tfrac12\log(1-x)+o(1)$.  Dominated convergence applies to $J$:
if $a=\arccos(\sqrt x\sin\theta)$, the absolute value of its integrand is
$a/\sin a\le\pi/2$.  Therefore
\[
 \lim_{x\uparrow1}J(x)
 =-\int_0^{\pi/2}\frac{u}{\sin u}\,du
 =2\int_0^{\pi/4}\log(\tan v)\,dv=-2G.
\]
It follows that
\begin{equation}\label{eq:Ffinitepart}
 \lim_{x\uparrow1}\left\{F(x)+\frac\pi{16}\log(1-x)\right\}
 =\frac\pi4\log2-\frac G2.
\end{equation}
Stirling's formula gives $t_j=\pi/(16j)+O(j^{-2})$.  Hence the series obtained by
subtracting $\pi/(16j)$ is absolutely convergent, and Abel's theorem turns
\eqref{eq:Ffinitepart} into \eqref{eq:Grecipseries}.  Rearrangement gives
\eqref{eq:Grecipseries2}.

Finally, the even recurrence in \eqref{eq:Trec}, or directly
\eqref{eq:Tsquares}, gives
\[
 \frac{T_{2m}}{\binom{2m}{m}^2}=-\sum_{j=1}^m t_j.
\]
Since $H_m-\log m\to\gamma$, the partial-sum form of
\eqref{eq:Grecipseries} proves \eqref{eq:Tevenfinitepart}.
\end{proof}

The endpoint $c=4$ is part of a full transport family.  It also explains why ordinary
binomial acceleration alone cannot turn this companion into an irrationality proof.

\begin{proposition}[binomial transport and its optimal acceleration]\label{prop:bintransport}
For a sequence $U$, set
\[
 U_n^{[c]}=\sum_{k=0}^n\binom nk(-c)^{n-k}U_k.
\]
If $U=A_E$, respectively $U=B_E$, then
\begin{align}\label{eq:transportrec}
 &(n+1)^2U_{n+1}^{[c]}
 -(4-c)(3n^2+3n+1)U_n^{[c]}
 +(32-24c+3c^2)n^2U_{n-1}^{[c]}\notag\\
 &\hspace{25mm}
 +c(4-c)(8-c)n(n-1)U_{n-2}^{[c]}
 =\begin{cases}0,&U=A_E,\\(-c)^n,&U=B_E.\end{cases}
\end{align}
Its characteristic polynomial factors as
\begin{equation}\label{eq:transportroots}
 (X-(8-c))(X-(4-c))(X+c).
\end{equation}
For $c=0,1,2,3$,
\begin{equation}\label{eq:transportlimit}
 \lim_{n\to\infty}\frac{B_n^{[c]}}{A_n^{[c]}}=\frac G2,\qquad
 \limsup_{n\to\infty}
 \left|\,\frac G2-\frac{B_n^{[c]}}{A_n^{[c]}}\,\right|^{1/n}
 =\frac{\max(c,4-c)}{8-c}.
\end{equation}
Thus $c=2$ is the unique optimum on $[0,4)$, improving the convergence factor from
$1/2$ to $1/3$.  Nevertheless no integral transform in this family produces a
linear form tending to zero: the numerator
$A_n^{[c]}G/2-B_n^{[c]}$ has exponential scale
$\max(c,4-c)\ge2$ for $c=0,1,2,3$, while $c=4$ has scale $4$ and the logarithmic
collision displayed in \eqref{eq:Tevenfinitepart}.
\end{proposition}

\begin{proof}
Ordinary binomial convolution gives
\[
 U^{[c]}(z)=\frac1{1+cz}U\!\left(\frac z{1+cz}\right).
\]
Substitute this identity into \eqref{eq:BOGF}; the chain rule gives
\eqref{eq:transportrec}.  The homogeneous leading part factors as
\eqref{eq:transportroots}.

For completeness, the nonvanishing of the relevant asymptotic components can be read
without a connection calculation.  The bounds \eqref{eq:Alower} and
$A_E(n)\le8^n$ give $A_E(n)^{1/n}\to8$.  The positive error formula
\eqref{eq:errorseries}, together with the recurrence, gives
\[
 \lim_{n\to\infty}|A_E(n)G/2-B_E(n)|^{1/n}=4.
\]
Under the generating-function substitution, these two singular scales move to
$8-c$ and $4-c$.  The latter singularity cannot disappear because the prefactor is
regular and nonzero at its preimage.  In addition, the right side of
\eqref{eq:transportrec} forces the simple characteristic root $-c$ in the transformed
linear form.  Poincar\'e asymptotics therefore give \eqref{eq:transportlimit}.
Finally,
\[
 \frac{\max(c,4-c)}{8-c}=
 \begin{cases}(4-c)/(8-c),&0\le c\le2,\\
 c/(8-c),&2\le c<4,\end{cases}
\]
which decreases to $1/3$ and then increases.  The asserted raw scales and the endpoint
case follow at once.
\end{proof}

\begin{theorem}[half-digit twisted Lucas law]\label{thm:CatalanLucas}
Let $p$ be an odd prime and let $0\le a,r<p$ with $2a+1<p$.  In the localisation
$\Zp$,
\begin{equation}\label{eq:CatalanLucas}
 p^2B_E(ap+r)\equiv
 \chi_{-4}(p)B_E(a)A_E(r)\pmod p.
\end{equation}
\end{theorem}

\begin{proof}
Write $k=bp+s$, $0\le s<p$, and extend the shell by zero outside $0\le k\le n$.
Lucas' theorem, applied also to the two central binomial coefficients, gives
\[
 S_E(ap+r,bp+s)\equiv S_E(a,b)S_E(r,s)\pmod p.
\]
Since $2a+1<p$, every harmonic argument in a contributing cell is strictly below
$p^2$.  Thus multiplying a weight-$q$ letter by $p^q$ makes it $p$-integral.

If $S_E(r,s)\not\equiv0\pmod p$, then Kummer's carry criterion gives
\[
 s\le r,\qquad 2s<p,\qquad 2(r-s)<p.
\]
Consequently
\[
 \left\lfloor\frac{k}{p}\right\rfloor=b,\qquad
 \left\lfloor\frac{2k}{p}\right\rfloor=2b,\qquad
 \left\lfloor\frac{2(n-k)}p\right\rfloor=2(a-b).
\]
Every monomial in the weight in \eqref{eq:CatalanB} has total weight $2$ and character
signature $\chi_{-4}$.  Lemma~\ref{lem:monomial} therefore gives
\[
 p^2W(ap+r,bp+s)\equiv\chi_{-4}(p)W(a,b)\pmod p.
\]
If $S_E(r,s)\equiv0$, the shell congruence and $p$-integrality of $p^2W$ make both
sides of the corresponding summand congruence zero.  Summing over the full digit block
separates the $b$- and $s$-sums and proves \eqref{eq:CatalanLucas}.
\end{proof}

\begin{lemma}[the last digit]\label{lem:CatalanLastDigit}
For every odd prime $p$,
\begin{equation}\label{eq:CatalanLastDigit}
 A_E(p-1)\equiv(-1)^{(p-1)/2}=\chi_{-4}(p)\pmod p.
\end{equation}
\end{lemma}

\begin{proof}
Put $h=(p-1)/2$.  In the sum defining $A_E(p-1)$, the factor
$\binom{2k}{k}$ vanishes modulo $p$ when $k>h$, while
$\binom{2p-2-2k}{p-1-k}$ vanishes when $k<h$.  Only $k=h$ remains.  Since
$\binom{p-1}{h}\equiv(-1)^h\pmod p$, its shell value is
$(-1)^{3h}=(-1)^h$.  The last equality in \eqref{eq:CatalanLastDigit} is Euler's
criterion for the quadratic character modulo $4$.
\end{proof}

\begin{theorem}[unit-digit full Lucas law]\label{thm:CatalanUnitLucas}
Let $p$ be an odd prime such that
\begin{equation}\label{eq:unitdigit}
 p\nmid A_E(j)\qquad(0\le j<p).
\end{equation}
Then for every $0\le a,r<p$ the unrestricted congruence
\begin{equation}\label{eq:CatalanUnitLucas}
 p^2B_E(ap+r)\equiv
 \chi_{-4}(p)B_E(a)A_E(r)\pmod p
\end{equation}
holds in $\Zp$.
\end{theorem}

\begin{proof}
Write $N=ap+r$.  Summing the Casoratian increments in the proof of
Corollary~\ref{cor:Catalanerror}, but stopping at $N$, gives the rational identity
\begin{equation}\label{eq:finitevariation}
 \frac{B_E(N)}{A_E(N)}=
 \sum_{m=1}^{N}\frac{32^{m-1}}{m^2A_E(m)A_E(m-1)}.
\end{equation}
Summing the shell factorisation used in the proof of
Theorem~\ref{thm:CatalanLucas} gives
$A_E(cp+s)\equiv A_E(c)A_E(s)\pmod p$.  Under \eqref{eq:unitdigit}, every
$A_E(m)$ with $0\le m<p^2$ is therefore a $p$-unit.  After multiplying
\eqref{eq:finitevariation} by $p^2$, all summands with $p\nmid m$ vanish modulo
$p$.  Writing the remaining indices as $m=jp$, Fermat's theorem and Lucas descent give
\begin{align*}
 \frac{p^2B_E(N)}{A_E(N)}
 &\equiv \sum_{j=1}^{a}
 \frac{32^{j-1}}{j^2A_E(j)A_E(j-1)A_E(p-1)}\\
 &=\frac{B_E(a)}{A_E(a)A_E(p-1)}\pmod p.
\end{align*}
Finally $A_E(N)\equiv A_E(a)A_E(r)$, and
$A_E(p-1)^{-1}\equiv\chi_{-4}(p)$ by
Lemma~\ref{lem:CatalanLastDigit}.  This proves \eqref{eq:CatalanUnitLucas}.
\end{proof}

\begin{corollary}[optimality of the denominator exponent]\label{cor:denoptimal}
For every prime $p\ge5$,
\begin{equation}\label{eq:vpBp}
 v_p(B_E(p))=-2.
\end{equation}
Consequently the exponent $2$ in Theorem~\ref{thm:sharpden} cannot be lowered in any
uniform bound $L_n^eB_E(n)\in\mathbb Z$.
\end{corollary}

\begin{proof}
Take $a=1$ and $r=0$ in Theorem~\ref{thm:CatalanLucas}; its half-digit condition is
$3<p$.  Since $B_E(1)=A_E(0)=1$, it gives
$p^2B_E(p)\equiv\chi_{-4}(p)\not\equiv0\pmod p$.  Theorem~\ref{thm:sharpden}
already gives $v_p(B_E(p))\ge-2$, so equality follows.
\end{proof}

\begin{remark}[the full-digit carry problem]\label{rem:fullLucas}
Exact recurrence computation gives \eqref{eq:CatalanLucas} for every
$0\le a,r<p$, without the half-digit restriction, for all odd primes $p\le31$.
The unrestricted statement does not follow termwise from
Theorem~\ref{thm:lucas}: when $2k$ crosses $p^2$, individual carry cells develop
$p$-adic poles.  Their sum nevertheless cancels in every test.  Proving this
second-digit carry cancellation would remove the residual zero-digit condition not
covered by Theorem~\ref{thm:CatalanUnitLucas} and yield a full companion Lucas law.
\end{remark}

\begin{remark}[irrationality boundary]
The formula is a new finite harmonic representation of the normalized second solution,
but not an irrationality proof for $G$.  The characteristic roots of
\eqref{eq:Erec-intro} are $8$ and $4$; the complementary linear form grows on the
$4^n$ scale instead of tending to zero, and clearing harmonic denominators only worsens
the arithmetic estimate.  This is consistent with the obstruction in earlier
Ap\'ery-like constructions for Catalan's constant \cite{ZudilinCatalanRemarks}.
Theorem~\ref{thm:sharpden} improves the clearing factor from the elementary
$L_{2n}^2=\exp((4+o(1))n)$ bound to $L_n^2=\exp((2+o(1))n)$, but the resulting
linear form still has size on the $4^nL_n^2$ scale.  Thus the new sharp denominator
theorem does not by itself imply irrationality.  Proposition~\ref{prop:bintransport}
sharpens this statement: even the optimal fixed integral binomial acceleration has a
raw small-form scale $2^n$, before any denominator is cleared.
\end{remark}

\section{Algorithm and further directions}

The results suggest the following reproducible search-and-proof pipeline.
\begin{enumerate}[leftmargin=2em]
\item Fit a candidate weight in a graded alphabet of ordinary and residue-coloured
  harmonic letters.  Separate character signatures from the beginning.
\item Convert every fitted monomial to its connected lift by
  Theorem~\ref{thm:universal}.  At weight two use Corollary~\ref{cor:weighttwo}; compact
  several blocks by matching logarithmic derivatives when possible.
\item Split indices only by the least common period needed for cyclotomic letters.
  Theorem~\ref{thm:cyclolift} then produces proper hypergeometric cells.
\item Telescope one-parameter polarizations such as \eqref{eq:polarization}, differentiate
  their recurrences, and eliminate lower jets.  Verify the rational certificates before
  using any initial values.
\item Test Theorem~\ref{thm:lucas} termwise.  If digit compatibility fails, the obstruction
  identifies exactly which wide argument or carry correction must be repaired.
\end{enumerate}

Two directions look particularly worthwhile.

\paragraph{Higher characters.}
The residue-indicator construction treats non-real characters without fractional
hypergeometric exponents.  It should make weight-two companions for cubic characters as
accessible as the $\chi_{-4}$ example, at the cost of a larger finite residue split.

\paragraph{Minimal operators from connected lifts.}
The universal lift proves existence but can produce a large Ore operator.  Connected
lifts contain much more cancellation than a generic parameter derivative.  A systematic
criterion for when those cancellations force the original Ap\'ery operator, rather than
a left multiple, would turn many experimentally minimal harmonic formulas into short
proofs.

\paragraph{Multi-digit divided-power jets.}
The single-digit bound in Theorem~\ref{thm:lucas} is honest.  A natural replacement for
iteration is to package the parameter Taylor series in divided powers, where denominators
and successive Frobenius descents can be tracked simultaneously.  This may connect the
elementary harmonic-jet picture with unit-root and crystal formulations of Ap\'ery
congruences.

\section{Reproducibility}

The accompanying script \nolinkurl{work/harmonic_jets/check_lifts.py} uses only Python's
integer and rational arithmetic.  It checks the logarithmic-derivative decomposition in
Proposition~\ref{prop:Catalanlift} cell by cell for $n\le19$ and audits
\eqref{eq:CatalanB} and \eqref{eq:Erec-intro} through $n=30$.  The script
\nolinkurl{work/harmonic_jets/verify_catalan_scalar.py} proves the scalar input
recurrences symbolically from \eqref{eq:increments}.  Finally
\nolinkurl{work/harmonic_jets/prove_catalan_ore.py} reconstructs the common factor,
verifies
the exact right divisions, and performs the finite propagation checks in
Proposition~\ref{prop:OreCatalan}; it runs in the open-source
PassageMath/\texttt{ore\_algebra} environment.  The Maxima file
\nolinkurl{work/harmonic_jets/explore_e.mac} records the parity-split setup, with
two-dimensional display disabled.  The exact audit
\nolinkurl{work/harmonic_jets/check_catalan_lucas.py} checks the unrestricted
congruence in Remark~\ref{rem:fullLucas} for every odd prime through $31$ and reports
which of them satisfy the proved unit-digit hypothesis.
The script \nolinkurl{work/harmonic_jets/check_catalan_endpoint_transform.py} checks
\eqref{eq:transportrec} for $0\le c\le8$, \eqref{eq:Trec}, the square expansion
\eqref{eq:Tsquares}, and the termwise divisibility $L_nR_{n,j}\in\mathbb Z$ through
$n=300$.  The SymPy script
\nolinkurl{work/harmonic_jets/verify_catalan_transport_symbolic.py} verifies the
generating-function chain rule and characteristic factorisation for symbolic $c$.
The independent script
\nolinkurl{work/harmonic_jets/verify_general_endpoint.py} verifies
\eqref{eq:bTrec} for symbolic $b$, audits \eqref{eq:bTsquares} through $n=120$ for
several positive and negative parameters divisible by $4$, and derives both necessity
witnesses in \eqref{eq:btestvalues} symbolically.
The older independent recurrence audit
\nolinkurl{work/harmonic_jets/check_catalan_denominators.py} verifies
Theorem~\ref{thm:sharpden} directly through $n=500$.
A detailed Lean handoff is provided in
\nolinkurl{work/harmonic_jets/LEAN_CATALAN_HANDOFF.md}.

\begin{thebibliography}{99}

\bibitem{BradleyCatalan}
D.~M.~Bradley,
\emph{A class of series acceleration formulae for Catalan's constant},
Ramanujan J. \textbf{3} (1999), no.~2, 159--173;
\texttt{arXiv:0706.0356}.

\bibitem{ABS2011}
J.~Ablinger, J.~Bl\"umlein and C.~Schneider,
\emph{Harmonic sums and polylogarithms generated by cyclotomic polynomials},
J. Math. Phys. \textbf{52} (2011), 102301; \texttt{arXiv:1105.6063}.

\bibitem{ChamStraub}
M.~Chamberland and A.~Straub,
\emph{Ap\'ery limits: experiments and proofs},
Amer. Math. Monthly \textbf{128} (2021), no.~9, 811--824;
\texttt{arXiv:2011.03400}.

\bibitem{KalmykovWardYost}
M.~Yu.~Kalmykov, B.~F.~L.~Ward and S.~A.~Yost,
\emph{Multiple (inverse) binomial sums of arbitrary weight and depth and the all-order
$\varepsilon$-expansion of generalized hypergeometric functions with one half-integer
value of parameter},
J. High Energy Phys. 2007, no.~10, 048; \texttt{arXiv:0707.3654}.

\bibitem{JiSunCLF}
X.-J.~Ji and Z.-H.~Sun,
\emph{Congruences for Catalan--Larcombe--French numbers},
Publ. Math. Debrecen \textbf{90} (2017), no.~3--4, 387--406;
\texttt{arXiv:1505.00668}.

\bibitem{MalikStraub}
A.~Malik and A.~Straub,
\emph{Divisibility properties of sporadic Ap\'ery-like numbers},
\texttt{arXiv:1508.00297}.

\bibitem{PWZ}
M.~Petkov\v{s}ek, H.~S.~Wilf and D.~Zeilberger,
\emph{$A=B$}, A K Peters, Wellesley, MA, 1996.

\bibitem{Phadikar}
J.~Phadikar,
\emph{Monoidal alphabets for generalized harmonic sums},
preprint (2026), \texttt{arXiv:2605.21525}.

\bibitem{RowlandYassawiKratt}
E.~Rowland, R.~Yassawi and C.~Krattenthaler,
\emph{Lucas congruences for the Ap\'ery numbers modulo $p^2$},
Integers \textbf{21} (2021), A20; \texttt{arXiv:2005.04801}.

\bibitem{Schneider2016}
C.~Schneider,
\emph{A difference ring theory for symbolic summation},
J. Symbolic Comput. \textbf{72} (2016), 82--127.

\bibitem{StraubZudilin}
A.~Straub and W.~Zudilin,
\emph{Sums of powers of binomials, their Ap\'ery limits, and Franel's suspicions},
Int. Math. Res. Not. \textbf{2023}, no.~11, 9861--9879;
\texttt{arXiv:2112.09576}.

\bibitem{StraubGessel}
A.~Straub,
\emph{Gessel--Lucas congruences for sporadic sequences},
Monatsh. Math. \textbf{203} (2024), 883--898;
\texttt{arXiv:2301.12248}.

\bibitem{WZ1990}
H.~S.~Wilf and D.~Zeilberger,
\emph{Rational functions certify combinatorial identities},
J. Amer. Math. Soc. \textbf{3} (1990), no.~1, 147--158.

\bibitem{WeiGong}
C.~Wei and D.~Gong,
\emph{Telescoping method, derivative operators and harmonic number identities},
Ramanujan J. \textbf{31} (2013), 177--188; \texttt{arXiv:1203.2051}.

\bibitem{Zagier}
D.~Zagier,
\emph{Integral solutions of Ap\'ery-like recurrence equations},
in: Groups and Symmetries, CRM Proc. Lecture Notes \textbf{47},
Amer. Math. Soc., 2009, 349--366.

\bibitem{ZudilinCatalanRemarks}
W.~Zudilin,
\emph{A few remarks on linear forms involving Catalan's constant},
\texttt{arXiv:math/0210423}.

\end{thebibliography}

\end{document}
